\section{Conclusion}
\label{sec:conclusion}

This paper introduced a strictly diagnostic framing of \emph{structural diagnosability} for enterprise-scale systems under incomplete observability.
The central contribution is not a new control mechanism, optimization strategy, or governance prescription, but a formal re-articulation of when
a system can be meaningfully diagnosed at all.

The first contribution is the formulation of \emph{diagnosability} as a property of the observable state space rather than of intervention capacity.
By explicitly separating observable manifestations from internal structure, the paper defines diagnosability in terms of whether the current observations
permit a non-empty and non-degenerate set of internally consistent structural hypotheses.
This reframing makes explicit that loss of diagnosability is a structural condition, not a failure of measurement effort or analytical sophistication.

The second contribution is the introduction of a phase-logic over diagnosability.
Rather than treating diagnosability as a continuous or monotonic quantity, the paper formalizes discrete diagnostic phases
(including stable diagnosis, marginal diagnosability, and non-diagnosable states) and the threshold transitions between them.
Within this logic, phase transitions correspond to structural changes in the admissible hypothesis space, not to managerial decisions or policy shifts.

The third contribution is the elevation of \emph{STOP} as a first-class diagnostic outcome.
STOP is not interpreted as inaction or failure, but as a formally valid conclusion reached when diagnosability collapses.
In such states, any prescriptive or optimization-oriented action would be epistemically unjustified.
Recognizing STOP as a legitimate terminal diagnostic result is essential for maintaining scientific coherence under severe uncertainty.

The falsifiability boundaries of the framework are explicit.
The core claims would be refuted if it could be shown that:
(a) diagnosability can be preserved despite an empty or internally inconsistent hypothesis space;
(b) phase transitions can be eliminated in favor of a fully continuous diagnostic regime without loss of explanatory power; or
(c) STOP-like terminal conditions can be replaced by universally valid prescriptive rules without additional structural assumptions.
Any such result would contradict the phase-logic and boundary conditions formalized here.

Equally explicit are the non-claims.
This paper does not propose methods for controlling, optimizing, steering, or stabilizing enterprises.
It does not provide decision rules, best practices, governance models, or performance improvements.
No claim is made that diagnosability implies controllability, nor that correct diagnosis enables successful intervention.
All prescriptive interpretations lie strictly outside the scope of this work.

This document depends on earlier formal artifacts for its definitions and primitives.
In particular, the ontological grounding and operator definitions are inherited from the archival artifacts commonly referred to as PDF-0 and PDF-1,
which provide the foundational formal context in which the present diagnostic analysis is situated.
The current paper should be read as a dependent, non-expansive contribution within that artifact lineage.

Future work is limited to artifact-level extensions rather than conceptual expansion.
This includes consolidation and expansion of the bibliographic base, formal tightening of selected proof sketches,
and the possible inclusion of executable or mechanically checkable demonstrations within the existing formal framework (e.g.\ as part of PDF-1–level artifacts).
No new diagnostic claims or prescriptive extensions are implied by these directions.

In summary, the paper establishes diagnosability as a structurally bounded property,
introduces a phase-logic for reasoning about its loss,
and legitimizes STOP as a scientifically necessary outcome.
Its value lies in clarifying limits rather than promising solutions,
and in restoring epistemic discipline where interventionist narratives are structurally unsupported.
